\documentclass[9.5pt,a4paper,landscape]{article}

\usepackage[utf8]{inputenc}
\usepackage[T1]{fontenc}
\usepackage[french]{babel}
\usepackage[a4paper,landscape,margin=0.8cm]{geometry}
\usepackage{lmodern}
\usepackage{xcolor}
\usepackage{tcolorbox}
\usepackage{array}
\usepackage{amsmath}

\pagestyle{empty}
\setlength{\parindent}{0pt}

\definecolor{bleu}{HTML}{245B78}
\definecolor{gris}{HTML}{F7F7F7}
\definecolor{corrige}{HTML}{1F6E43}

\newtcolorbox{exo}[1]{
  colback=white,
  colframe=black!55,
  boxrule=0.45pt,
  arc=1.2mm,
  left=1.8mm,
  right=1.8mm,
  top=1.5mm,
  bottom=1.3mm,
  height=5.35cm,
  valign=top,
  title=\textbf{Exercice #1},
  coltitle=bleu,
  fonttitle=\small,
  colbacktitle=gris,
  titlerule=0.25pt
}

\newcommand{\rep}[1]{%
  {\color{corrige}\bfseries #1}%
}

\begin{document}

\fontsize{9}{10.5}\selectfont

\begin{center}
{\large\bfseries\color{bleu}
Séquence 1 -- Statistiques : corrigé des exercices hors classe}\par
\end{center}

\vspace{0.25em}

%==================================================
% LIGNE 1
%==================================================

\noindent
\begin{minipage}[t]{0.242\linewidth}
\begin{exo}{1}
Voitures déjà comptées :

\[
8+12+14=34
\]

Il reste :

\[
50-34=16
\]

\rep{Il y a 16 voitures rouges.}
\end{exo}
\end{minipage}
\hfill
%
\begin{minipage}[t]{0.242\linewidth}
\begin{exo}{2}
La fréquence est :

\[
\frac{6}{24}
=
\frac14
\]

Or :

\[
\frac14=25\,\%
\]

\rep{La fréquence est \(\frac14\), soit \(25\,\%\).}
\end{exo}
\end{minipage}
\hfill
%
\begin{minipage}[t]{0.242\linewidth}
\begin{exo}{3}
On cherche \(25\,\%\) de 32 :

\[
32\times\frac{25}{100}=8
\]

\rep{8 élèves ont choisi l'option théâtre.}
\end{exo}
\end{minipage}
\hfill
%
\begin{minipage}[t]{0.242\linewidth}
\begin{exo}{4}
\[
\frac{8+11+9+12+10}{5}
=
\frac{50}{5}
=
10
\]

\rep{La moyenne est 10.}
\end{exo}
\end{minipage}

\vspace{0.16cm}

%==================================================
% LIGNE 2
%==================================================

\noindent
\begin{minipage}[t]{0.242\linewidth}
\begin{exo}{5}
Nombre total de scores :

\[
2+3+1=6
\]

\[
\frac{2\times8+3\times10+1\times12}{6}
=
\frac{58}{6}
\approx 9{,}7
\]

\rep{Le score moyen est environ \(9{,}7\).}
\end{exo}
\end{minipage}
\hfill
%
\begin{minipage}[t]{0.242\linewidth}
\begin{exo}{6}
On range les valeurs dans l'ordre croissant :

\[
11;\ 12;\ 14;\ 17;\ 19;\ 26;\ 30
\]

Il y a 7 valeurs.

La valeur du milieu est la 4\ieme :

\[
\rep{\text{Médiane}=17}
\]
\end{exo}
\end{minipage}
\hfill
%
\begin{minipage}[t]{0.242\linewidth}
\begin{exo}{7}
On range les valeurs dans l'ordre croissant :

\[
13;\ 15;\ 18;\ 19;\ 21;\ 27
\]

Il y a 6 valeurs.

Les deux valeurs centrales sont \(18\) et \(19\).

\[
\frac{18+19}{2}=18{,}5
\]

\[
\rep{\text{Médiane}=18{,}5}
\]
\end{exo}
\end{minipage}
\hfill
%
\begin{minipage}[t]{0.242\linewidth}
\begin{exo}{8}
Au moins la moitié des valeurs sont

\[
\rep{\text{inférieures ou égales à }14}
\]

et au moins la moitié sont

\[
\rep{\text{supérieures ou égales à }14}.
\]
\end{exo}
\end{minipage}

\vspace{0.16cm}

%==================================================
% LIGNE 3
%==================================================

\noindent
\begin{minipage}[t]{0.242\linewidth}
\begin{exo}{9}
On range les valeurs dans l'ordre croissant :

\[
-96;\ -92;\ -87;\ -86;\ -82
\]

Il y a 5 valeurs.

La valeur du milieu est la 3\ieme :

\[
\rep{\text{Médiane}=-87\text{ m}}
\]
\end{exo}
\end{minipage}
\hfill
%
\begin{minipage}[t]{0.242\linewidth}
\begin{exo}{10}
\textbf{Population :}

\rep{les 28 élèves interrogés.}

\medskip

\textbf{Caractère étudié :}

\rep{le nombre de frères et sœurs de chaque élève.}
\end{exo}
\end{minipage}
\hfill
%
\begin{minipage}[t]{0.242\linewidth}
\begin{exo}{11}
Sous forme de fraction :

\[
\frac{14}{40}
=
\frac{7}{20}
\]

En écriture décimale :

\[
\frac{14}{40}=0{,}35
\]

En pourcentage :

\[
0{,}35=35\,\%
\]

\rep{La fréquence est \(35\,\%\).}
\end{exo}
\end{minipage}
\hfill
%
\begin{minipage}[t]{0.242\linewidth}
\begin{exo}{12}
Une interprétation possible :

\medskip

\rep{Si la durée totale des trajets était répartie également entre tous les élèves, chacun aurait un trajet de 18 minutes.}
\end{exo}
\end{minipage}

\end{document}